\section{Proof of the main theorem}\label{sec:proof}

This section chooses the parameters and adds up the bounds of \cref{sec:accounting}. Throughout, $k\in\N$, $n=2k$ is the number
of heads, and $\eps>0$.

\subsection{The parameters}\label{sec:proof-params}

\begin{definition}[The parameters]\label{def:params}
\formal{Main.lean}{parH, parQ, parJ, parP, parK}
The \emph{grid width}~$h_0$, the \emph{bit size}~$Q$, the number of \emph{blocks}~$j$, the number of \emph{digits}~$P$ and the
\emph{radix}~$K$ are
\[
\begin{aligned}
  h_0&=\Bigl\lceil\frac{14n+80n^2}{\eps/4}\Bigr\rceil+1, &\qquad
  Q&=\Bigl\lceil\frac{2n^2}{\eps/4}\Bigr\rceil+1, &\qquad
  j&=\Bigl\lceil\ln\bigl(\max(1,8n^2/\eps)\bigr)\cdot Q^{n^2}\Bigr\rceil,\\
  P&=1+2n(jn+1), &
  K&=\Bigl\lceil\frac{5Pn+140\cdot2^{h_0}n^2}{\eps/4}\Bigr\rceil+1. &&
\end{aligned}
\]
\end{definition}

The parameters depend only on $k$ and~$\eps$, and they are fixed in this order: $h_0$ and~$Q$ from $n$ and~$\eps$, $j$ from~$Q$,
$P$ from~$j$, and $K$ from $P$ and~$h_0$. By construction $h_0,Q,K\ge1$ and $P\ge1$. The grid width pays for the terms that decay
like $1/h_0$, the bit size for exposure, the number of blocks, and with it the number of digits, for truncation, and the radix for
the wraps. The cut will have depth at most $K^P-1$.

\begin{lemma}[The parameters suffice]\label{lem:params}
\formal{Main.lean}{parH_spec, parQ_spec, parJ_spec, parK_spec}
\formal{Choice.lean}{le_natCeil_add_one_mul, one_sub_pow_le_inv, one_sub_pow_natCeil_le_inv, trunc_term_le}
The parameters of \cref{def:params} satisfy
\[
  14n+80n^2\le h_0\frac\eps4,\qquad 2n^2\le Q\frac\eps4,\qquad 2n^2\bigl(1-Q^{-n^2}\bigr)^j\le\frac\eps4,\qquad
  5Pn+140\cdot2^{h_0}n^2\le K\frac\eps4.
\]
\end{lemma}

\begin{outline}
For $\delta>0$ and every real~$C$, $C\le(\lceil C/\delta\rceil+1)\delta$, where $\lceil\cdot\rceil$ is the ceiling truncated at~$0$
(\lean{le_natCeil_add_one_mul}). With $\delta=\eps/4$ this gives the first, the second and the
fourth inequality (\lean{parH_spec}, \lean{parQ_spec}, \lean{parK_spec}). For the third (\lean{parJ_spec}, \lean{trunc_term_le}),
put $x=Q^{-n^2}\in(0,1]$ and $A=\max(1,8n^2/\eps)\ge1$, so that $j=\lceil\ln A/x\rceil$. Then $(1-x)^j\le e^{-xj}\le1/A$
(\lean{one_sub_pow_le_inv}, \lean{one_sub_pow_natCeil_le_inv}), and $2n^2/A\le\eps/4$: if $8n^2/\eps\le1$, then $2n^2\le\eps/4$,
and otherwise $2n^2/A=\eps/4$.
\end{outline}

\subsection{The cut, summed over the random choices}\label{sec:proof-total}

As in \cref{sec:accounting}, $\Omega=[0,h_0)\times[0,K)^P\times[0,Q)^\iota$ is the set of the random choices $(r,\omega,\xi)$,
$\Psi_{r,\omega,\xi}$ is the label of \cref{def:global}, and $\mathrm{Bad}_{r,\xi}$ is the set of bad times of \cref{lem:trunc}.
The following proposition adds up the bounds of \cref{sec:accounting} in natural numbers, for arbitrary parameters.

\begin{proposition}[The cut in natural numbers]\label{prop:total}
\formal{Main.lean}{sum_weight_events_head_le, sum_violated_head_le, sum_labelCutSize_le, sum_bad_le}
Let $h_0,P\in\N$ and $K,Q\ge1$.
\begin{enumerate}[label=(\alph*)]
\item For every head~$v$, $\sum_{r<h_0}\sum_{(t,g)}2^{\theta_r(g)}\le8T$, the inner sum over
$\mathrm{Rise}_r(v)\cup\mathrm{Last}_r(v)$.
\item For every head~$v$, the upper hops $(p,u)$ of~$v$ with $\Psi_{r,\omega,\xi}(u)\le\Psi_{r,\omega,\xi}(p)$, summed over
$(r,\omega,\xi)\in\Omega$, number at most
\begin{multline*}
  B_v=K^PQ^\iota\cdot6T+K^P\sum_{r<h_0}\sum_\xi\card{\mathrm{Bad}_{r,\xi}}+K^PQ^\iota\cdot80nT\\
  +K^Ph_0T\cdot2nQ^{\iota-1}+h_0Q^\iota K^{P-1}\cdot20\cdot2^{h_0}n(6T+1).
\end{multline*}
\item The cuts, summed over~$\Omega$, satisfy
\[
  \sum_{(r,\omega,\xi)\in\Omega}\card{F_{\Psi_{r,\omega,\xi}}}\le K^PQ^\iota\cdot8nT+Q^\iota h_0\cdot PK^{P-1}\cdot n(4T+1)+n\cdot B,
\]
where $B$ is the common value of the bounds~$B_v$ of~(b).
\item If $P\ge1+2n(jn+1)$ for some $j\in\N$, then $\sum_{r<h_0}\sum_\xi\card{\mathrm{Bad}_{r,\xi}}\le h_0T\cdot2n(1-Q^{-n^2})^jQ^\iota$.
\end{enumerate}
\end{proposition}

\begin{outline}
(a) \Cref{lem:budgets}(c) with $N=T$, and \cref{lem:W1}(b) (\lean{sum_weight_events_head_le}). (b) (\lean{sum_violated_head_le})
By \cref{lem:hopterms}(a), for each choice the violated upper hops of~$v$ are at most its lost, truncated, charged, exposed and
wrapping hops; we sum each kind over~$\Omega$, the offsets in the sums being below~$h_0$. \emph{Lost:} the lost set depends only
on~$r$, and by Theorem~A (\cref{thm:A}) its size, summed over the offsets, is at most $6$ times the number of upper hops, hence at
most $6T$ (\cref{lem:totals}(a)); the other two coordinates contribute the factor $K^PQ^\iota$. \emph{Truncated:} at most
$\card{\mathrm{Bad}_{r,\xi}}$ for every~$\omega$ (\cref{lem:trunc}(a), \lean{card_truncSet_le_bad}). \emph{Charged:} the charged
set depends only on~$r$; by \cref{lem:hopterms}(b) and part~(a), summed over the offsets, its size is at most $10\sum_w8T=80nT$.
\emph{Exposed:} independent of~$\omega$, and summed over~$\xi$ at most $T\cdot2nQ^{\iota-1}$ for every~$r$ (\cref{lem:hopterms}(c),
\cref{lem:totals}(a)). \emph{Wrapping:} for every $r$ and~$\xi$, summed over~$\omega$, at most
$K^{P-1}\cdot20\cdot2^{h_0}\sum_w(6T+1)$ (\cref{lem:hopterms}(d), \cref{lem:totals}(e)). (c) (\lean{sum_labelCutSize_le}) The cut
consists of the descents and the violated hop edges. \Cref{lem:heads}(c) bounds the latter by the violated upper hops of the
$n$ heads, each of which is bounded by~(b), and \cref{lem:totals}(f) bounds the descents. (d) (\lean{sum_bad_le}) Exchanging the
sums, $\sum_\xi\card{\mathrm{Bad}_{r,\xi}}=\sum_{t<T}\#\{\xi:t\text{ is bad}\}$, and \cref{lem:trunc}(b) bounds every term.
\end{outline}

\subsection{The main estimate}\label{sec:proof-main}

The four error terms of the main estimate are, as expectations over uniform random choices, the following. The terms \emph{paid
by the grid} are the descents at event times, $8nT/h_0$, the lost hops, $6nT/h_0$ (Theorem~A), and the charged hops,
$80n^2T/h_0$; the term of \emph{truncation} is $2n^2(1-Q^{-n^2})^jT$ (Lemma~D); the term of \emph{exposure} is $2n^2T/Q$
(\cref{prop:exposure}); and the term of the \emph{wraps} collects the wrap times, $Pn(4T+1)/K$, and the wrapping hops,
$20\cdot2^{h_0}n^2(6T+1)/K$. With the parameters of \cref{def:params}, each term is at most $\eps T/4$.

\begin{proposition}[The main estimate]\label{prop:mainest}
\formal{Main.lean}{exists_avgLabel}
Let $T\in\N$ and let $x_1,\dots,x_k$ be walks, pauses allowed. With the parameters of \cref{def:params}, there are a nonempty
finite set~$A$ and labels $\Psi_a$, $a\in A$, with values below~$K^P$, such that
\[
  \sum_{a\in A}\card{F_{\Psi_a}}\le\card A\cdot\eps T .
\]
\end{proposition}

\begin{outline}
The formal proof takes $A=\Omega$ and the labels $\Psi_{r,\omega,\xi}$ of \cref{def:global}. They take fewer than $K^P$ values by
\cref{lem:heads}(a), and $\Omega$ is nonempty. For $T=0$ there are no edges. Let
$T\ge1$. Since $P=1+2n(jn+1)$, \cref{prop:total}(d) applies; insert it into \cref{prop:total}(c) and divide by
$\card\Omega=h_0K^PQ^\iota$. The result is a sum of four terms, each at most $\eps T/4$ by \cref{lem:params}:
\begin{enumerate}[label=(\roman*)]
\item \emph{paid by the grid:} $(8n+6n+80n^2)T/h_0=(14n+80n^2)T/h_0\le\eps T/4$;
\item \emph{truncation:} $n\cdot2n(1-Q^{-n^2})^jT=2n^2(1-Q^{-n^2})^jT\le\eps T/4$;
\item \emph{exposure:} $n\cdot2nTQ^{\iota-1}/Q^\iota=2n^2T/Q\le\eps T/4$; here $Q^\iota=Q\cdot Q^{\iota-1}$, since for $k\ge1$ every
head has the root as a frame, so that $\iota\ge1$, and for $k=0$ the term vanishes;
\item \emph{wraps:} $n\bigl(P(4T+1)+20\cdot2^{h_0}n(6T+1)\bigr)/K\le(5Pn+140\cdot2^{h_0}n^2)T/K\le\eps T/4$, since $4T+1\le5T$ and
$6T+1\le7T$.
\end{enumerate}
So the average cut is at most $\eps T$ (\lean{exists_avgLabel}).
\end{outline}

\subsection{Theorem T}\label{sec:proof-T}

\begin{theorem}[Theorem T]\label{thm:T}
\formal{Main.lean}{tcpAt_direct, tcp, main, sr}
For $k\in\N$ and $\eps>0$ let $D(k,\eps)=K^P-1$, with $K$ and~$P$ as in \cref{def:params}. For every horizon~$T$ and all walks
$x_1,\dots,x_k$, with or without pauses, the walk DAG of $x_1,\dots,x_k$ before~$T$ has a cut of at most $\eps T$ edges and depth
at most $D(k,\eps)$. Consequently:
\begin{enumerate}[label=(\alph*)]
\item $\TCP(k)$ holds for every~$k$, for walks with pauses, and hence for pause-free walks: this is the main statement;
\item $\SR(s)$ holds for every~$s$: DAGs with a spine whose arcs form $s$ stacks reduce to constant depth.
\end{enumerate}
\end{theorem}

\begin{outline}
The first statement is \lean{tcpAt_direct}, that is, $\TCP(k)$ at $(\eps,D(k,\eps))$ in the sense of \cref{def:tcp}
(\lean{BoundedNesting.TCPat}). By \cref{prop:mainest} some nonempty finite family of labels with values below~$K^P$ (in the formal
proof, the labels $\Psi_{r,\omega,\xi}$) has cuts of average size at most $\eps T$; by the probabilistic method (\cref{prop:avg},
\lean{exists_cut_of_avg}) one of these labels has a cut of at most $\eps T$ edges, whose depth is at most $K^P-1$ by
\cref{prop:labelcut}. (a) For every $\eps>0$ the first statement provides
the depth $D(k,\eps)$ (\lean{tcp}); the main statement, $\TCP(k)$ for every~$k$ (\lean{MainStatement}), is \lean{main}. Walks
without pauses are walks, so $\TCP(k)$ contains its pause-free form (\lean{BoundedNesting.tcpnp_of_tcp}). (b) (\lean{sr},
\lean{sr_of_main}) By the reduction of the eighth paper (\cref{prop:stacksred}, \lean{BoundedNesting.sr_of_tcpnp}), an $s$-stack
DAG on $[0,N)$ is realized by $s$ pause-free walks of length~$2N$, and a cut of their walk DAG of at most $(\eps/2)(2N)$ edges and
depth~$D$ pulls back to a cut of the stack DAG of at most $\eps N$ edges and depth at most~$D$. The formal proof takes the depth
of~(a) at $\eps/2$, so every $s$-stack DAG on $[0,N)$ has a cut of at most $\eps N$ edges and depth at most $D(s,\eps/2)$.
\end{outline}

The construction never uses pause-freeness: Theorem~A and the composition hold for walks with pauses, so the first statement covers
them directly, at the depth of $k$ walks, without the reduction to $2k$ pause-free walks of \cref{prop:stacksred}. The depth
$D(k,\eps)$ does not depend on~$T$, so no family of $k$-tuples of walks is depth-robust, and with~(b) every family of DAGs with a
spine whose arcs form a bounded number of stacks is depth-reducible; in particular $\SR(3)$, the first open case left by the
eighth paper~\cite[Problem~10.2]{EighthPaper2026}, holds.
